\documentclass{beamer}

\usetheme{Copenhagen}
\setbeamertemplate{frametitle}[default][center]

\title{Comparative Analysis of Machine Learning Approaches\\for MNIST Digit Recognition}
\author{CheesyChocolate}
\date{\today}

\begin{document}

\begin{frame}
\maketitle
\end{frame}

\section{Introduction}

\begin{frame}
\frametitle{Introduction}

\begin{itemize}
    \item The MNIST dataset is a fundamental benchmark in computer vision
    \item 70,000 handwritten digits (60,000 for training, 10,000 for testing)
    \item Each image is a 28$\times$28 grayscale pixel image
    \item 10 classes (digits 0-9)
\end{itemize}

\begin{figure}
    \centering
    \includegraphics[width=0.7\textwidth]{fig/mnist_samples.png}
    \caption{Sample MNIST digits}
\end{figure}

\end{frame}

\section{Project Goals}

\begin{frame}
\frametitle{Project Goals}

\begin{itemize}
    \item Implement multiple pattern recognition approaches for MNIST classification
    \item Compare and contrast different algorithms:
    \begin{itemize}
        \item Traditional ML: SVM, Random Forest, KNN
        \item Deep Learning: CNN, MLP, CycleGAN, Vision Transformer
    \end{itemize}
    \item Analyze performance metrics across methods
    \item Explore semi-supervised learning for limited labeled data scenarios
    \item Identify strengths and weaknesses of each approach
\end{itemize}

\end{frame}

\section{Methodology}

\begin{frame}
\frametitle{Methodology}

\begin{itemize}
    \item Data preprocessing
    \begin{itemize}
        \item Normalization (0-255 → 0-1)
        \item Used a subset of 5,000 images for computational efficiency
        \item 80\% training, 20\% testing split
    \end{itemize}
    \item Standardized evaluation metrics
    \begin{itemize}
        \item Accuracy, Precision, Recall, F1 Score
        \item Confusion matrices
    \end{itemize}
\end{itemize}

\end{frame}

\section{Models}

\begin{frame}
\frametitle{Traditional Machine Learning Approaches}

\begin{itemize}
\item \textbf{Support Vector Machine}
    \begin{itemize}
        \item RBF kernel
        \item C=10.0, gamma='scale'
        \item Finds optimal decision boundaries
    \end{itemize}
\item \textbf{Random Forest}
    \begin{itemize}
        \item Ensemble of 200 decision trees
        \item Bootstrapped training
        \item Majority voting
    \end{itemize}
\item \textbf{K-Nearest Neighbors}
    \begin{itemize}
        \item k=3 neighbors
        \item Uniform weighting
        \item Instance-based learning
    \end{itemize}
\end{itemize}

\end{frame}

\begin{frame}
\frametitle{Deep Learning Approaches}

\begin{itemize}
\item \textbf{Convolutional Neural Network (CNN)}
    \begin{itemize}
        \item Two convolutional layers (32 and 64 filters)
        \item Max pooling and dropout for regularization
        \item Fully connected layer with 128 neurons
    \end{itemize}
\item \textbf{Multi-Layer Perceptron (MLP)}
    \begin{itemize}
        \item Two hidden layers (512 and 256 neurons)
        \item ReLU activation functions
        \item Dropout regularization
    \end{itemize}
\item \textbf{Vision Transformer (ViT)}
    \begin{itemize}
        \item Divides images into 7×7 patches
        \item 4 transformer layers with 4 attention heads
        \item Optimized for CPU execution
        \item Adds position embeddings to capture spatial information
    \end{itemize}
\end{itemize}

\end{frame}

\begin{frame}
\frametitle{Semi-Supervised Learning with CycleGAN}

\begin{itemize}
\item \textbf{CycleGAN for Limited Labeled Data}
    \begin{itemize}
        \item Adversarial architecture with generators and discriminators
        \item Leverages unlabeled data through feature learning
        \item Adds classification head to discriminator for digit recognition
    \end{itemize}
\item \textbf{Semi-Supervised Approach}
    \begin{itemize}
        \item Uses only 500 labeled examples (vs. 60,000 in full dataset)
        \item Two-phase training: GAN training + classifier fine-tuning
        \item Demonstrates effective learning when labels are scarce
    \end{itemize}
\item \textbf{Real-World Applications}
    \begin{itemize}
        \item Medical imaging with limited annotated examples
        \item Industrial inspection with few labeled defects
        \item Document analysis with sparse labeled samples
    \end{itemize}
\end{itemize}

\end{frame}

\section{Results}

\begin{frame}
\frametitle{Accuracy Comparison}

\begin{figure}
    \centering
    \includegraphics[width=0.8\textwidth]{fig/model_comparison.png}
    \caption{Model accuracy comparison}
\end{figure}

\end{frame}

\begin{frame}
\frametitle{Confusion Matrices - Traditional Models}

\begin{figure}
    \centering
    \includegraphics[width=0.45\textwidth]{fig/SVM_confusion_matrix.png}
    \includegraphics[width=0.45\textwidth]{fig/RandomForest_confusion_matrix.png}
    \caption{SVM (left) and Random Forest (right) confusion matrices}
\end{figure}

\end{frame}

\begin{frame}
\frametitle{More Confusion Matrices - Traditional and Transformer}

\begin{figure}
    \centering
    \includegraphics[width=0.45\textwidth]{fig/KNN_confusion_matrix.png}
    \includegraphics[width=0.45\textwidth]{fig/Transformer_confusion_matrix.png}
    \caption{KNN (left) and Vision Transformer (right) confusion matrices}
\end{figure}

\end{frame}

\begin{frame}
\frametitle{Confusion Matrices - Neural Networks}

\begin{figure}
    \centering
    \includegraphics[width=0.45\textwidth]{fig/CNN_confusion_matrix.png}
    \includegraphics[width=0.45\textwidth]{fig/MLP_confusion_matrix.png}
    \caption{CNN (left) and MLP (right) confusion matrices}
\end{figure}

\end{frame}

\begin{frame}
\frametitle{Sample Predictions - Traditional Models}

\begin{figure}
    \centering
    \includegraphics[width=0.95\textwidth]{fig/SVM_sample_predictions.png}
    \caption{Sample predictions from SVM model}
\end{figure}

\end{frame}

\begin{frame}
\frametitle{Sample Predictions - Neural Networks}

\begin{figure}
    \centering
    \includegraphics[width=0.95\textwidth]{fig/CNN_sample_predictions.png}
    \caption{Sample predictions from CNN model}
\end{figure}

\end{frame}

\begin{frame}
\frametitle{CycleGAN Training and Results}

\begin{figure}
    \centering
    \includegraphics[width=0.45\textwidth]{fig/cyclegan_discriminator_loss.png}
    \includegraphics[width=0.45\textwidth]{fig/cyclegan_classifier_training.png}
    \caption{CycleGAN training: Discriminator loss (left) and classifier accuracy (right)}
\end{figure}

\end{frame}

\begin{frame}
\frametitle{CycleGAN Sample Predictions}

\begin{figure}
    \centering
    \includegraphics[width=0.95\textwidth]{fig/CycleGAN_sample_predictions.png}
    \caption{Sample predictions from CycleGAN model with limited labeled data}
\end{figure}

\end{frame}

\section{Conclusions}

\begin{frame}
\frametitle{Conclusions}

\begin{itemize}
    \item CNN achieves the highest accuracy at 96.00\%
    \item SVM follows at 95.40\%
    \item Random Forest achieves 93.00\%
    \item KNN performs at 92.80\%
    \item MLP and Transformer both at 89.50\%
    \item CycleGAN demonstrates effective semi-supervised learning with limited labeled data
    \item Trade-offs between accuracy, speed, and interpretability
    \begin{itemize}
        \item CNN: Best accuracy but requires more computational resources
        \item SVM: Good balance of accuracy and efficiency
        \item Random Forest: Good interpretability
        \item KNN: Simple baseline but memory-intensive
        \item MLP: Decent performance with simpler architecture than CNN
        \item CycleGAN: Effective when labeled data is scarce
        \item Transformer: Competitive performance on CPU without specialized hardware
    \end{itemize}
\end{itemize}

\end{frame}

\begin{frame}
\frametitle{Hardware Considerations}

\begin{itemize}
    \item Vision Transformer was successfully optimized for CPU usage
    \begin{itemize}
        \item Reduced model complexity (fewer layers, smaller batch size)
        \item Added explicit AMD GPU (ROCm) support detection
        \item Optimized memory usage with pin\_memory
        \item Properly initialized position embeddings
    \end{itemize}
    \item Shows modern architectures can be adapted for resource-constrained environments
    \item Performance trade-offs are acceptable for many applications
    \item Democratizes access to advanced deep learning models
\end{itemize}

\end{frame}

\begin{frame}
\frametitle{Future Work}

\begin{itemize}
    \item Further hyperparameter optimization for each model
    \item Feature engineering to enhance digit representation
    \item Exploration of recurrent and hybrid architectures
    \item Ensemble methods combining multiple models
    \item Scaling up the Vision Transformer with more data
    \item Optimizing models for mobile and edge devices
    \item Application to more complex datasets
\end{itemize}

\end{frame}

\begin{frame}
\frametitle{Thank You!}

\begin{center}
\huge{Questions?}
\end{center}

\end{frame}

\end{document}
